\documentclass[11pt]{article}

\usepackage{amsmath, amssymb, amsthm}
\usepackage{graphicx}
\usepackage{hyperref}
\usepackage{tikz}
\usepackage{enumitem}
\usepackage{geometry}
\usepackage{algorithm}
\usepackage{algpseudocode}
\usepackage[backend=biber,
    maxnames=10,
    minnames=1,
style=alphabetic]{biblatex}
\addbibresource{references.bib}
\geometry{margin=1in}

\newtheorem{proposition}{Proposition}
\newtheorem{theorem}{Theorem}
\newtheorem{lemma}{Lemma}
\newtheorem{corollary}{Corollary}

\title{Integer Programming of Graph Squares}
\author{Eddie Xu\\\small\texttt{eddie.xu@student.unimelb.com.au}}
\date{\today}


\begin{document}

\maketitle

\begin{abstract}
In this paper, we investigate the matrix properties of graph squares, focusing in particular on the
relationship between the adjacency matrix of a graph and the corresponding adjacency of its square. 
\end{abstract}

\tableofcontents

\section{Introduction}

The problem of finding the square root of a graph is a well-studied field in both mathematics and computer-science, 
with specific relation to distributed computing, network design and computational biology. Research in this area has
primarily focused around the undirected graph $H$ and its square $G = H^2$ which is obtained by adding an edge between 
every pair of vertices whose distance in $H$ is at most two. 

Much of the existing research on graph square roots has centered on characterizing when a graph admits such 
a root and on the complexity of reconstructing $H$ from $G$. In contrast rather than studying the construction problem 
directly, we analyze the relationship between the adjacency matrix $A_H$ of the root and the adjacency matrix $A_G$ of 
the square, and show that the problem can be completely categorized as an integer program. 

\section{Background}

Let $H = (V, E_H)$ be the undirected \emph{graph root} with $|V|=n$ vertices, and let $G = (V, E_G)$ be the undirected 
\emph{graph square} where the relationship between root $H$ and square $G$ is denoted as $G = H^2$. The
\emph{adjacency matrix} of $G$ can be represented as $A_G$ where edge $ij \in E_G$ exists if and only if $i \neq j$ and 
there is a path in $H$ of length $1$ or $2$ from $i$ to $j$. Equivalently, 
\[
(A_G)_{ij} =
\begin{cases}
1 ,& \text{if } \big( (A_H)_{ij} = 1 \;\text{or}\; (A_H^2)_{ij} > 0 \big) \; \text{for} \; i \neq j, \\
0 ,& \text{otherwise}.
\end{cases}
\]

This definition follows strictly from $(A_H^2)_{ij}$ giving the number of paths of length two between 
vertices $i, j$. We however will use an 
alternative definition with the \emph{degree matrix} $(D_H)_{ii} = \deg_H(i)$ such that 
\[
A_{G = H^2} = \frac{(A_H)^2 - D_H}{\# \; \text{paths between} \; i, j} + A_H. 
\]
Let $\mathbf{1}_{>0} : \mathbb{R}^{n \times n} \to \{0,1\}^{n \times n}$ denote the entrywise \emph{indicator function}, defined by
\[
\mathbf{1}_{>0}(M)_{ij} =
\begin{cases}
1 & \text{if } M_{ij} > 0, \\
0 & \text{if } M_{ij} = 0.
\end{cases}
\]
Then the adjacency matrix of the graph square $G$ is given by 
\begin{equation} \label{eq::fund_adj_eq_undirected}
    A_G = \mathbf{1}_{> 0}\!\big((A_H)^2 - D_H + A_H\big).
\end{equation}

\section{Integer Program}

Given that an edge exists in the graph if 

\begin{proposition}
    Let $a,b,c \in V$ and let $ab, bc \in E_G, ca \notin E_G$. Then 
    \begin{enumerate}
        \item $(A_H)_{ab} = 1$, $(A_H)_{bc} = 0$ with $d \in V$ such that , or
         $(A_H)_{ab} = 0$, $(A_H)_{bc} = 1$, or 
        \item $(A_H)_{ab} = (A_H)_{bc} = 0$. 
    \end{enumerate}
\end{proposition}

\begin{proof}
    If $(A_H)_{ab} = 1$, $(A_H)_{bc} = 1$, then edge $ca$ must exist in the square. 
\end{proof}

\begin{proposition}
    Let $a,b,c \in V$ and let $ab, bc, ca \in E_G$. Then 
    \begin{enumerate}
        \item $(A_H)_{ab} = (A_H)_{bc} = (A_H)_{ca} = 1$, or
        \item $(A_H)_{ab} = 0$, $(A_H)_{bc} = (A_H)_{ca} = 1$ or $(A_H)_{bc} = 0$, $(A_H)_{ab} = (A_H)_{ca} = 1$ or 
        $(A_H)_{ca} = 0$, $(A_H)_{ab} = (A_H)_{bc} = 1$, or
        \item $(A_H)_{ab} = (A_H)_{bc} = (A_H)_{ca} =  0$. 
    \end{enumerate}
\end{proposition}

\begin{proof}
    If $(A_H)_{ab} = 1$, $(A_H)_{bc} = 1$, then edge $ca$ must exist in the square. 
\end{proof}


\printbibliography

\end{document}
